% ============================================================================
% 04_mockups.tex — Cuerpo del entregable "Mockups"
% Wireframes de baja/media fidelidad (PlantUML Salt), sin colores de marca.
% Sin preámbulo: se incluye tanto desde 04_Mockups.tex (standalone) como,
% eventualmente, desde un documento maestro.
% ============================================================================
\section{Mockups}

\subsection{Propósito y alcance}

Esta sección presenta los \textbf{mockups} (bocetos de interfaz) de las
cinco pantallas de LSA Salud, en fidelidad baja a media, sin colores de
marca ni tipografía definitiva: el objetivo es comunicar la disposición de
la información y las interacciones disponibles, no el diseño visual final
(que sí existe en la implementación real, \texttt{web/app/**/page.tsx} y
\texttt{web/components/**}, con estilos de Tailwind CSS, para las cuatro
pantallas ya implementadas). Los wireframes se generaron con la notación
\textbf{Salt} de PlantUML a partir de una lectura directa de esos
componentes, de modo que cada elemento mostrado corresponde a un elemento
real de la interfaz implementada; la excepción es P-05 (Pantalla de
Consentimiento), diseño nuevo sin componente real todavía (RF-03 oficial).

Cada elemento relevante de una pantalla está marcado con un \textbf{callout}
numerado entre paréntesis --(1), (2), (3)...--, único dentro de esa pantalla.
Estos callouts son referenciados más adelante por el Entregable~07 (Casos de
Uso Reales) para describir la interacción concreta sobre la interfaz real,
por ejemplo: ``el sistema muestra la palabra reconocida en el panel (6) de
la Pantalla P-02''.

Las cinco pantallas identificadas son las mismas de la sección 3.1 de la
ERS (ver Entregable 05) y del registro de identificadores (\texttt{IDS.md}):
P-01 (Selección de Estación), P-02 (Estación del Paciente), P-03 (Estación
de Recepción), P-04 (Modo Captura) y P-05 (Pantalla de Consentimiento,
diseñada, pendiente de implementación).

Un principio de producto atraviesa las pantallas del flujo de comunicación y
se remarca en cada sección: en el flujo del paciente \textbf{no hay botones
de confirmación manual}. La oración traducida se exhibe brevemente en
pantalla (ventana de exhibición previa, RF-01.6) y ese lapso actúa como
confirmación automática antes de reproducir el audio; el paciente no tiene
que pulsar ``aceptar'' ni ``enviar'' (ver RNF-T-05). La única confirmación
manual explícita del flujo es, precisamente, la nueva Pantalla P-05 de
consentimiento (RF-03.1): es una excepción deliberada al principio, porque
el consentimiento de cámara exige una acción positiva y consciente del
paciente, no una confirmación automática.

% ============================================================================
\subsection{P-01 --- Selección de Estación}

Landing page servida en la ruta \texttt{/} (\texttt{web/app/page.tsx}).
Permite abrir la Estación del Paciente o la Estación de Recepción, pensada
para un entorno de doble monitor o dos dispositivos.

\diagrama{04_mockup_p01_landing}{0.85\textwidth}{Mockup de P-01 --- Selección de Estación}{fig:mock-p01}

\begin{table}[H]
    \centering
    \caption{Elementos de la Pantalla P-01}
    \label{tab:mock-p01}
    \begin{tabularx}{\textwidth}{c X}
        \toprule
        \textbf{Callout} & \textbf{Elemento} \\
        \midrule
        (1) & Logotipo de LSA Salud. \\
        (2) & Título de bienvenida (``Bienvenido a LSA Salud''). \\
        (3) & Texto explicativo: instrucción de abrir cada estación en una
        ventana o pantalla distinta. \\
        (4) & Tarjeta ``Estación del Paciente''; enlaza a P-02
        (\texttt{/paciente}). \\
        (5) & Tarjeta ``Estación de Recepción''; enlaza a P-03
        (\texttt{/recepcion}). \\
        \bottomrule
    \end{tabularx}
\end{table}

% ============================================================================
\subsection{P-02 --- Estación del Paciente (modo predicción)}

Pantalla servida en \texttt{/paciente} (\texttt{web/app/paciente/page.tsx})
cuando el modo activo es \texttt{predict}. Concentra la cámara, el
reconocimiento de señas en tiempo real y la conversación compartida con la
Estación de Recepción.

\diagrama{04_mockup_p02_paciente}{\textwidth}{Mockup de P-02 --- Estación del Paciente}{fig:mock-p02}

\begin{table}[H]
    \centering
    \caption{Elementos de la Pantalla P-02}
    \label{tab:mock-p02}
    \begin{tabularx}{\textwidth}{c X}
        \toprule
        \textbf{Callout} & \textbf{Elemento} \\
        \midrule
        (1) & Encabezado (\texttt{Header}) con el logotipo de LSA Salud. \\
        (2) & Título de sección y descripción de la estación. \\
        (3) & Indicador ``EN VIVO · Visión por computadora'' (estado activo
        de la cámara; solo alcanzable tras el consentimiento de P-05, RF-03.1). \\
        (4) & Área de video con el flujo de la cámara y los landmarks
        de MediaPipe superpuestos en un \texttt{canvas}. \\
        (5) & Ventana de oración en construcción, superpuesta al video:
        muestra las palabras reconocidas mientras la secuencia sigue
        abierta, o el indicador ``Generando respuesta de IA...'' mientras
        se traduce (RF-01.4, RF-01.5). Se retira sola cuando la oración se
        cierra o se traduce: no requiere confirmación del paciente. \\
        (6) & Panel ``Top-3 Predicciones'' (\texttt{TopPredictions}): tres
        barras con la seña y el porcentaje de confianza de las tres clases
        más probables en cada ciclo de predicción (RF-01.3). \\
        (7) & Panel de conversación: lista de burbujas de chat
        (\texttt{ChatBubble}) del paciente y del personal de recepción o
        enfermería, sincronizada entre estaciones (RF-01.7). \\
        (8) & Botón ``Limpiar Historial''. \\
        (9) & Botón ``Ajustes (Modo Captura)''; cambia a la Pantalla P-04. \\
        (10) & Ícono ``Reportar'' sobre cada burbuja del panel (7); abre
        CU-10 (Reportar Traducción Incorrecta o Ambigua, RF-06.1).
        \textbf{Diseñado --- pendiente de implementación.} \\
        \bottomrule
    \end{tabularx}
\end{table}

Nótese que no hay ningún botón de confirmación entre (5) y (7): la oración
pasa de ``en construcción'' a integrar la conversación de forma automática
en cuanto el sistema cierra y traduce la secuencia (principio de producto,
RNF-T-05). El callout (10) es la única excepción deliberada al principio de
``sin confirmación manual'': reportar una traducción es, por definición,
una acción explícita del paciente o del personal.

% ============================================================================
\subsection{P-03 --- Estación de Recepción}

Pantalla servida en \texttt{/recepcion} (\texttt{web/app/recepcion/page.tsx}).
Muestra el historial de la conversación y permite responder por texto o por
dictado de voz.

\diagrama{04_mockup_p03_recepcion}{\textwidth}{Mockup de P-03 --- Estación de Recepción}{fig:mock-p03}

\begin{table}[H]
    \centering
    \caption{Elementos de la Pantalla P-03}
    \label{tab:mock-p03}
    \begin{tabularx}{\textwidth}{c X}
        \toprule
        \textbf{Callout} & \textbf{Elemento} \\
        \midrule
        (1) & Encabezado (\texttt{Header}) con el logotipo de LSA Salud. \\
        (2) & Título ``Estación de Recepción / Enfermería'' y descripción. \\
        (3) & Panel ``Historial del Paciente'': burbujas de chat
        (\texttt{ChatBubble}) de la conversación compartida (RF-01.7). \\
        (4) & Área de texto para redactar una respuesta
        (\texttt{ReceptionControls}). \\
        (5) & Botón ``Enviar''; habilitado solo con texto no vacío (RF-02.1). \\
        (6) & Botón circular de dictado por voz (micrófono); alterna entre
        escuchar y detener (RF-02.2, CU-05). \\
        (7) & Indicador de estado: ``Click para hablar'' / ``Escuchando...
        click para detener'', o el mensaje de error si el navegador no
        soporta reconocimiento de voz o falla el micrófono. \\
        (8) & Ícono ``Reportar'' sobre cada burbuja del panel (3); abre
        CU-10 (Reportar Traducción Incorrecta o Ambigua, RF-06.1).
        \textbf{Diseñado --- pendiente de implementación.} \\
        \bottomrule
    \end{tabularx}
\end{table}

% ============================================================================
\subsection{P-04 --- Modo Captura}

Sub-pantalla de \texttt{/paciente} (\texttt{web/components/capture/CaptureMode.tsx})
activada al elegir ``Ajustes (Modo Captura)'' desde P-02. Permite grabar
nuevas muestras de entrenamiento para el modelo de reconocimiento.

\diagrama{04_mockup_p04_captura}{0.95\textwidth}{Mockup de P-04 --- Modo Captura}{fig:mock-p04}

\begin{table}[H]
    \centering
    \caption{Elementos de la Pantalla P-04}
    \label{tab:mock-p04}
    \begin{tabularx}{\textwidth}{c X}
        \toprule
        \textbf{Callout} & \textbf{Elemento} \\
        \midrule
        (1) & Logotipo de LSA Salud, en el encabezado del Modo Captura. \\
        (2) & Selector de etiqueta de seña: chips con el nombre de cada
        seña conocida y, entre paréntesis, la cantidad de muestras ya
        grabadas para esa etiqueta (RF-T-02); la etiqueta seleccionada
        queda resaltada. \\
        (3) & Botón ``+ Nueva seña''; despliega un campo de texto para
        crear una etiqueta que todavía no existe (RF-T-01, CU-07). El
        nombre se valida contra el patrón \texttt{\^{}[a-z0-9\_]\{1,32\}\$}
        antes de aceptarse. \\
        (4) & Buscador de señas (filtro de texto); aparece solo cuando hay
        más de seis etiquetas conocidas, para facilitar encontrar una en
        una lista larga. \\
        (5) & Vista previa de la cámara con los landmarks superpuestos,
        igual que en P-02. \\
        (6) & Cuenta regresiva (3, 2, 1) antes de empezar a grabar cada
        repetición (RF-T-03). \\
        (7) & Indicador ``GRABANDO'', visible mientras se capturan los 45
        frames de la muestra actual. \\
        (8) & Indicador ``MUESTRA \emph{n} / \emph{N}''; solo se muestra
        cuando se configuraron más de una repetición. \\
        (9) & Control numérico ``Repeticiones''; cantidad de muestras a
        grabar de corrido para la etiqueta seleccionada. \\
        (10) & Botón ``Empezar'' (inicia la secuencia de grabación) que se
        reemplaza por ``Detener'' mientras hay una grabación en curso
        (extensión 6a de CU-06: detener cancela la repetición actual sin
        guardarla). \\
        (11) & Botón ``Volver a Inferencia''; regresa a P-02. \\
        \bottomrule
    \end{tabularx}
\end{table}

% ============================================================================
\subsection{P-05 --- Pantalla de Consentimiento}
\textit{Diseñado --- pendiente de implementación. Cubre el requisito oficial RF-03.}

Pantalla modal prevista para \texttt{/paciente}, mostrada antes de invocar
\texttt{navigator.mediaDevices.getUserMedia()} y antes de cualquier otro
elemento de P-02. Bloquea el acceso a la cámara hasta que el paciente
otorgue el consentimiento explícito (CU-09; RF-03 oficial). Diseño nuevo,
sin componente real todavía; se describe aquí en el mismo nivel de
fidelidad que el resto de los wireframes para dejar constancia de su
disposición y contenido mínimo.

\diagrama{04_mockup_p05_consentimiento}{0.75\textwidth}{Mockup de P-05 --- Pantalla de Consentimiento}{fig:mock-p05}

\begin{table}[H]
    \centering
    \caption{Elementos de la Pantalla P-05 (diseñada, pendiente de implementación)}
    \label{tab:mock-p05}
    \begin{tabularx}{\textwidth}{c X}
        \toprule
        \textbf{Callout} & \textbf{Elemento} \\
        \midrule
        (1) & Logotipo de LSA Salud, en el encabezado del modal. \\
        (2) & Título ``Consentimiento para usar la cámara''. \\
        (3) & Texto explicativo en lenguaje simple: qué capta la cámara,
        para qué se usa (traducir señas) y que no se guarda video ni
        imágenes (RNF-03 oficial). \\
        (4) & Botón ``Acepto y continúo''; otorga el consentimiento
        (CU-09) y activa la cámara. \\
        (5) & Botón ``No acepto''; deja la cámara inactiva y muestra una
        alternativa (por ejemplo, solicitar asistencia del personal). \\
        \bottomrule
    \end{tabularx}
\end{table}

Al igual que en P-02/P-03, el botón (4) es la única acción manual
requerida en todo el flujo del paciente: es una excepción deliberada al
principio de ``sin confirmación manual'' (ver introducción de esta
sección), porque el consentimiento exige una decisión explícita y
consciente, no una confirmación automática por temporización.
